%! Unit 1: Algebraic Expressions and Factoring — the unit packet cover sheet.
%
% THIS FILE IS THE COVER. Both wrappers \input it — unit_cover/main.tex (student
% packet) and unit_cover_key/main.tex (key packet, which adds a page 2 of exam
% scoring notes) — so page 1 can never drift between the two packets. Edit the
% cover here, never in a wrapper.

% Full-bleed royal banner
\begin{tikzpicture}[remember picture, overlay]
  \fill[royal] (current page.north west)
    rectangle ([yshift=-0.9in]current page.north east);
\end{tikzpicture}
\vspace{-0.8in}

\begin{center}
  {\color{white}\textbf{\LARGE Algebra, Trigonometry, and Data Analysis}}\\[4pt]
  {\color{mistmid}\large\textbf{Unit 1: Algebraic Expressions and Factoring}}\\[6pt]
  {\color{white}\textbf{\large Unit Overview \quad\textbar\quad Eight Lessons \quad\textbar\quad Unit Test}}
\end{center}

\vspace{0.20in}
\namedateperiod
\vspace{0.14in}

\noindent
\small
Everything in this unit is one habit run in two directions. \textbf{Multiplying} takes two
factors and hides them inside a single expression; \textbf{factoring} takes that expression
apart and puts the factors back. Every lesson here is a different shape of that same move ---
and once an expression is factored, two new questions open up that you could not answer before:
\emph{when is it zero?} (Lesson 1.6) and \emph{what cancels?} (Lesson 1.7).

\vspace{0.12in}

% BOXGUARD: all three boxes below are breakable and all three fit on page 1 today
% with ~1.5in to spare, so these guards are inert as written. They are sized to
% their own boxes so that if the cover ever grows, a box moves whole to a second
% page instead of splitting and stranding a stub — the guard is the cheap
% insurance, and each was verified not to cost this sheet its single page.
\boxguard[14]
\begin{tocbox}
\small
\renewcommand{\arraystretch}{1.32}
\begin{tabularx}{\linewidth}{@{} c >{\raggedright\arraybackslash}p{5.5cm} l X @{}}
\rowcolor{mist}
\textbf{\#} & \textbf{Lesson} & \textbf{SOL} & \textbf{The one idea} \\
\midrule
1.0 & Unit Opener                    & ---                 & One model, two equivalent forms \\
1.1 & Expressions \& Exponent Rules  & \texttt{A2.EO.3a}   & An exponent counts factors, not terms \\
1.2 & Radicals \& Rational Exponents & \texttt{A2.EO.2a--c}& A radical is an exponent rule run backwards \\
1.3 & Factoring: GCF \& Grouping     & \texttt{A2.EO.3b}   & Pull out the \emph{largest} factor, not the first \\
1.4 & Factoring Trinomials           & \texttt{A2.EO.3b}   & Split the middle term, then group \\
1.5 & Factoring Special Forms        & \texttt{A2.EO.3b, d}& Patterns you already earned --- and can verify \\
1.6 & Solving by Factoring           & \texttt{A2.EI.2b, 6a--b} & Zero is the only product worth having \\
1.7 & Rational Expressions           & \texttt{A2.EO.1a--b}& Only \emph{factors} cancel \\
\end{tabularx}
\end{tocbox}

\vspace{0.12in}

\boxguard[16]
\begin{spiralbox}
\small
\begin{enumerate}[leftmargin=1.5em, itemsep=3pt, topsep=2pt]
  \item \textbf{An operation does not reach inside a sum.} $(2m+3)^2 \ne 4m^2+9$,
        $\sqrt{9+16} \ne 3+4$, and $\frac{x+5}{x+7} \ne \frac57$. Three lessons, one belief ---
        squaring, rooting, and cancelling all distribute over \emph{multiplication} only.
  \item \textbf{``Completely'' means largest, not first.} $6x(4x+6)$ is a true statement and a
        wrong answer. Check every factor you write down to see whether it factors again.
  \item \textbf{Factoring recovers information.} Multiplying erases the dimensions of a rectangle;
        factoring an area puts them back.
  \item \textbf{Zero is the only number that tells you anything.} $ab=12$ says nothing about $a$
        or $b$; $ab=0$ forces one of them to be zero. That is why every equation gets moved to
        standard form first.
  \item \textbf{A cancelled factor leaves a scar.} Excluded values are read from the
        \emph{original} denominator. Simplifying hides a restriction; it never removes one.
\end{enumerate}
\end{spiralbox}

\vspace{0.12in}

\boxguard[8]
\begin{remindbox}
\small
\textbf{The unit test has four parts and is worth 100 points:} Part A vocabulary matching
(8 pts), Part B multiple choice (12 pts), Part C short answer and computation (50 pts), and
Part D extended response (30 pts). A \textbf{practice test} with the same structure and
different numbers is bound at the back of this packet --- work it with no notes, then check it
against the key. Part D asks you to \emph{justify}, not just to compute: expect to verify an
identity by multiplying, and to say where two expressions agree and where they do not.
\end{remindbox}
